% ECONOMICS_PROBLEM: {"id": "F31-399", "title": "Exchange-rate disconnect", "jel_code": "F31", "source_id": "OP-0331"}
% !TeX program = xelatex
\documentclass[11pt,a4paper]{article}
\usepackage[margin=23mm]{geometry}
\usepackage{fontspec}
\setmainfont{Latin Modern Roman}
\usepackage{amsmath,amssymb,mathtools}
\usepackage{xcolor,enumitem,booktabs,longtable,array,xurl,hyperref}
\definecolor{muted}{HTML}{53636C}
\hypersetup{colorlinks=true,urlcolor=blue,linkcolor=blue}
\setlength{\parindent}{0pt}
\setlength{\parskip}{5pt}
\newcommand{\R}{\mathbb R}
\newcommand{\E}{\mathbb E}
\newcommand{\N}{\mathbb N}
\newcommand{\ind}{\mathbf 1}
\newcommand{\Var}{\operatorname{Var}}
\newcommand{\Cov}{\operatorname{Cov}}
\newcommand{\supp}{\operatorname{supp}}
\DeclareMathOperator*{\argmin}{arg\,min}
\DeclareMathOperator*{\argmax}{arg\,max}
\newcommand{\entrymeta}[3]{{\sffamily\small\color{muted}\textbf{\texttt{#1}}\quad|\quad #2\par #3\par}}
\newcommand{\Problemref}[1]{\texttt{#1}}
\newcommand{\JELref}[2]{\texttt{#2} (JEL \texttt{#1})}
\begin{document}
\section*{Exchange-rate disconnect}
\textbf{Problem ID:} \texttt{F31-399}\quad \textbf{JEL:} \texttt{F31}\par
% BEGIN ECONOMICS STATEMENT
\label{problem:OP-0331}
{\sffamily\small\color{muted}Original subject group: International finance and exchange rates\par}
\label{definitions:problem:30007613}
\textbf{Audit classification:} Published research agenda. The historical source names the disconnect puzzle and motivates research. The 2025 manuscript provides a substantive productivity-expectations account. The four-group variance decomposition is editorial and requires a corrected stationary/cointegrated representation.
\subsection*{Definitions and precise research target}
Fix exchange-rate data \(s_t\), current and future productivity, consumption, trade, financing conditions and measured expectations for specified countries and dates. Let \(y_t\) collect a stationary transformation of these observables (with differencing, trend or cointegration treatment explicitly stated); use the stationary exchange-rate component as the row denoted by \(s\) below. A structural representation is
\[y_t=\sum_{h\geq0}B_h\varepsilon_{t-h},\qquad E[\varepsilon_t\varepsilon_t']=I,\]
where \(B_h\) are response matrices and \(\varepsilon_t\) are mean-zero structural shocks divided into current fundamentals, noisy news about future fundamentals, financial-demand shocks and intermediary-capacity shocks. The summation must converge in mean square; \(I\) is the identity covariance matrix. The task is to identify the exchange-rate forecast-error variance shares attributable to each group at declared horizon \(H\geq1\), restricting attention to positive exchange-rate forecast-error variance,
\[\theta_g(H)=\frac{\sum_{h=0}^{H-1}\sum_{k\in g}(B_h)_{s,k}^2}{\sum_{h=0}^{H-1}\sum_k(B_h)_{s,k}^2}.\]
Here \((B_h)_{s,k}\) is the response of \(s\) to shock \(k\) and group \(g\) is specified before estimation. Provide restrictions distinguishing news from contemporaneous fundamentals and financial shocks; sign restrictions alone may yield sets rather than unique shares. Examine whether apparent weak contemporaneous correlations reflect different response horizons rather than absence of fundamental links. Report identified sets and robustness to omitted shocks, exchange regimes and survey measurement errors. The decomposition requires an identification design and does not follow from low exchange-rate correlations alone.

\textbf{Additional formal definitions and scope (editorial).} The displayed moving-average representation is appropriate for covariance-stationary variables. Since nominal log exchange-rate levels are often integrated, specify \(y_t\) using differences, detrended stationary variables, or an explicit cointegrated representation; do not impose a square-summable stationary representation on an untreated unit-root level. Let \(\sum_h\|B_h\|_F^2<\infty\) and shocks be orthogonal across dates with identity within-date covariance. Forecast-error variance at horizon \(H\) must use information revealing past structural innovations, or justify invertibility from observed histories. If forecasting the level from innovations in \(\Delta s\), use accumulated level-response coefficients rather than the one-period differenced coefficients in the numerator. Orthogonal rotations preserve observable covariance, so identifying shock groups requires restrictions beyond the reduced form.

For the identification part, declare an admissible class \(\Theta\) of structural models, including preferences, technologies, shock laws, measurement/sampling rules and any equilibrium selector. If \(O\) denotes the complete observable history and \(T(\theta)\) the target defined above, the identified set is \(\mathcal I_T=\{T(\theta):\theta\in\Theta,\ \mathcal L_\theta(O)=\mathcal L_{\rm obs}(O)\}\); point identification means this set is a singleton. A \(do(a)\) comparison replaces stated policy/technology equations while fixing the declared exogenous law and re-solving the remaining equations. These definitions do not assert that an identification theorem holds without further restrictions.
\subsection*{Variants and assumption changes}
\begin{enumerate}
\item \textbf{Productivity-news explanation (source-motivated).} Impose restrictions separating anticipated productivity from noisy signals and quantify identified variance shares.
\item \textbf{Financial and intermediary alternatives (editorial).} Add funding/demand shocks and survey measurement, using a stationary or cointegrated representation; bound decompositions over observationally equivalent rotations.
\end{enumerate}
\subsection*{References and source audit}
\begin{enumerate}
\item Publisher verifies Obstfeld and Rogoff, volume 15 (2000), pp.339--390 and DOI; copyright 2001 is distinguished from volume year. Locator: Publisher abstract, second paragraph. Support: Historical disconnect puzzle. Abstract and metadata inspected; full article not independently inspected. URL: \url{https://www.journals.uchicago.edu/doi/10.1086/654423}.
\item Author-hosted cover verifies five authors and November 21, 2025 manuscript; NBER WP 32596 DOI is the working-paper identifier. Locator: Abstract and Introduction, printed pp.1--4. Support: Productivity-expectations explanation. Full manuscript accessible; version distinguished from initial WP issue. URL: \url{https://chahrour.github.io/papers/Disconnect_revisited.pdf}.
\end{enumerate}
\begin{enumerate}\raggedright
\item\hypertarget{definitions:source:30007613:1}{} Maurice Obstfeld; Kenneth Rogoff. 2000. \emph{The Six Major Puzzles in International Macroeconomics: Is There a Common Cause?}. Publisher abstract, second paragraph.
\url{https://www.journals.uchicago.edu/doi/10.1086/654423}.
DOI: \url{https://doi.org/10.1086/654423}.
\item\hypertarget{definitions:source:30007613:2}{} Ryan Chahrour; Vito Cormun; Pierre De Leo; Pablo Guerrón-Quintana; Rosen Valchev. 2025. \emph{Exchange Rate Disconnect Revisited}. Abstract and Introduction, printed pp.1--4.
\url{https://chahrour.github.io/papers/Disconnect_revisited.pdf}.
DOI: \url{https://doi.org/10.3386/w32596}.
\end{enumerate}

\subsection*{Common conventions: Source mathematical conventions}

These conventions apply to each entry unless explicitly replaced.

\textbf{Models and identification.} A primitive model \(m\) specifies all probability laws, feasible choices, information, equilibrium and measurement rules used by its target. The admissible class \(\mathcal M\) is fixed independently of outcomes. If \(P_O(m)\) is its observable probability law and \(\tau(m)\) the target, its identified set at observable law \(P\) is
\[
 \mathcal I_\tau(P)=\{\tau(m):m\in\mathcal M,\ P_O(m)=P\}.
\]
Point identification means that this set is a singleton. An empty set signals incompatibility of data and assumptions. Coordinate bounds need not describe the joint set. Bounds are sharp only if every retained target is attainable or arbitrarily approachable by compatible models. Finite bounds require explicit restrictions on unobserved quantities; unrestricted confounding, hidden wealth, extrapolation or arbitrary equilibrium selection can yield uninformative sets.

\textbf{Causal interventions.} A potential outcome \(Y_i(p)\) is the outcome under a complete policy regime \(p\), including timing, financing, information and market spillovers. The target population and units are fixed before treatment. With interference, use the complete assignment vector or a justified exposure mapping; individual no-interference notation alone cannot identify an economy-wide intervention. A difference of mean potential outcomes does not require the cross-world joint distribution. Conditioning on post-treatment migration, survival, enrollment or employment changes the estimand and needs separate justification.

\textbf{Statistical statements.} An estimator, confidence set, forecast or test specifies a sampling law, model class and loss/coverage criterion. Uniform \((1-\alpha)\)-coverage means \(\inf_{m\in\mathcal M}\Pr_m\{\tau(m)\in C_n\}\geq1-\alpha\), or an explicitly asymptotic version. The whole parameter space trivially has coverage, so informativeness requires a stated width, risk or power objective. A forecasting improvement is relative to a named benchmark, information set and evaluation population.

\textbf{Equilibrium and optimization.} An equilibrium comprises feasible plans, optimality, beliefs where needed, budget constraints and all market/resource clearing. The primitives or a stated benchmark define each correspondence \(\mathcal E(p;m)\). Nonemptiness is assumed or proved, never inferred just from compact policy sets. A selected-equilibrium target uses a declared selection rule; a robust target quantifies over all permitted equilibria. Use suprema or infima unless attainment is established. An \(\varepsilon\)-optimal policy is feasible and within \(\varepsilon\) of the finite optimum value. Report an empty feasible set separately.

\textbf{Welfare and accounting.} Normative utility, interpersonal normalization, social weights, numeraire, discounting and the use of public receipts are primitives. Behavioral utility need not coincide with the welfare criterion. Household welfare includes ownership income and rents once; adding profits again to compensating variation generally double-counts them. A monetary equivalent is defined by an explicit transfer and price convention, with monotonicity and range conditions. Average effects, mechanism contributions, transfers and welfare losses are different targets. Infinite sums require convergence and dynamic feasible plans require terminal or no-Ponzi conditions.

\textbf{Source versions.} A working-paper date, revision date and journal publication year are distinguished. A DOI for a journal version is not automatically the DOI of an earlier manuscript. Locators refer to the stated version and printed pages where specified. A metadata-only check does not verify a claim about a full-text conclusion. Every entry's source subsection narrows the provenance claim accordingly.
% END ECONOMICS STATEMENT
\end{document}
